% !TeX root = ../main.tex
% 章边界：给定外生道路恢复历程和参考配置规则，研究服务过程如何演化。
% 道路恢复行动的内生优化留给第七章。
\chapter{道路可达性与资源供需耦合下的服务韧性演化}
\label{chap:resilience}

\section{道路恢复与时变供需下的服务状态转移}
\subsection{外生恢复历程与通行能力变化}
\subsection{资源补给、需求到达与状态耦合}

\ThesisBlank

\section{需求时限与服务损失累积}
\subsection{资源等待、运输延迟与需求到期}
\subsection{阶段服务缺口与不可补偿损失}

\ThesisBlank

\section{服务韧性轨迹与演化阶段刻画}
\subsection{服务下降、受限运行与恢复过程}
\subsection{不同韧性轨迹的比较与分解}

\ThesisBlank

\section{服务恢复条件与恢复次序依赖分析}
\subsection{道路恢复与服务恢复的条件区分}
\subsection{需求到期与恢复次序的交互影响}

\ThesisBlank

\section{动态服务瓶颈辨识及其转移条件}
\subsection{道路、库存与时限约束的瓶颈判别}
\subsection{恢复阶段变化下的瓶颈转移}

\ThesisBlank

\section{服务韧性演化的机制实验与泸定情景分析}
\subsection{相同初末状态下不同恢复历程的配对比较}
\subsection{恢复历程与需求时序的交互检验}
\subsection{瓶颈变化与未恢复情景的反事实分析}
\subsection{不同恢复条件下的韧性差异与适用边界}

\ThesisBlank

\section{本章小结}
\ThesisBlank
